\section{Abstract}
In large cities, finding a parking space is an everyday problem. The steady growth of the vehicle fleet and the limited urban infrastructure mean that a significant share of traffic consists of vehicles circulating without a concrete destination in search of a free spot, increasing congestion, fuel consumption and polluting emissions. At the same time, a large number of private garages and parking lots remain empty for much of the day due to the lack of a channel that reliably connects them with demand.

This Capstone Project documents the design, development and implementation of \textit{ParkLink}, a digital platform that connects drivers who need to park with owners of private garages and parking lots, using the Autonomous City of Buenos Aires (CABA) as its initial case study. The solution consists of a mobile application for Android and iOS, with differentiated experiences for drivers and owners, and a web administration panel intended for the platform's staff. Drivers can search for parking spaces on a georeferenced map, check their real availability, book them by the hour or for longer stays, and pay from within the application. Owners, in turn, publish their spaces, set prices, schedules and capacity, approve incoming requests and receive the settlement of their earnings. The system also includes in-person check-in verification through codes, messaging and support tied to each booking, user reviews, push and real-time notifications, refund management and listing moderation.

From a technical standpoint, the platform was built on an architecture of independent microservices, each with its own database, accessed through a single API gateway responsible for validating authentication through JWT tokens and routing requests. The mobile application and the administration panel were developed with React Native and Expo. The backend follows a polyglot approach: Python with FastAPI for the gateway, authentication and notifications; Go for booking and availability management; NestJS with TypeScript for parking space management; and Node.js with Express for payments. Persistence relies on PostgreSQL, using the PostGIS extension for proximity-based geographic searches, and asynchronous communication between services is based on events published over RabbitMQ. External services were also integrated for payment processing (Mercado Pago), push notifications (Firebase Cloud Messaging), email and SMS verification (Brevo and Twilio), image storage (Cloudflare R2) and error monitoring (Sentry), with all components containerized using Docker.

Development was managed with a hybrid agile methodology based on Scrum and Kanban practices, organized into fixed-length iterations with planning and story-point estimation, and supported by periodic meetings with the tutors to validate each increment. The source code was organized in a GitHub organization with an independent repository per microservice, integrated into a root repository through submodules, with code review through pull requests and continuous integration with GitHub Actions running static analysis and automated tests.

As a result, a functional Minimum Viable Product (MVP) was obtained that validates the technical and operational feasibility of the proposal, leaving more complex integrations
 for later stages. ParkLink thus helps reduce parking search time and unnecessary circulation, and promotes a more efficient use of existing urban infrastructure.
